%! TEX root = ../main.tex

\clearpage

\section{Các yêu cầu phi chức năng}

\subsection{Yêu cầu thực thi}

\textit{<Nếu có các yêu cầu thực thi đối với sản phẩm cho các tình huống khác nhau, ta sẽ trình bày 
chúng trong phần này và giải thích lý do cơ bản để giúp nhà phát triển hiểu mục đích và tạo ra các 
lựa chọn thiết kế phù hợp. Với các hệ thống thời gian thực, ta phải xác định mối quan hệ theo thời 
gian. Các yêu cầu thực thi cần phải viết cụ thể như có thể. Ta có thể trình bày các yêu cầu thực thi 
theo từng tính năng hay từng yêu cầu chức năng riêng lẻ.> }

\subsection{Yêu cầu an toàn}

\textit{<Xác định các yêu cầu có liên quan tới các khả năng mất, hư hại hay hỏng do sử dụng sản phẩm. 
Định nghĩa các hoạt động hay bộ phận an toàn cần phải được thực hiện cũng như các hoạt động 
cần phải bị ngăn chặn. Tham khảo tới các quy tắc, các chính sách bên ngoài mà chúng trình bày 
các vấn đề về an toàn có ảnh hưởng đến việc thiết kế hay sử dụng sản phẩm. Nêu các giấy chứng 
nhận an toàn cần phải được đáp ứng.> }

\subsection{Yêu cầu bảo mật}
\textit{<Xác định các yêu cầu liên quan đến các vấn đề bảo mật và đời tư xung quanh việc sử dụng sản 
phẩm hay sự bảo vệ dữ liệu được sử dụng hay được tạo ra bởi sản phẩm. Định nghĩa các yêu cầu xác thực danh tính người sử dụng. Tham khảo tới các quy tắc, các chính sách bên ngoài có các 
vấn đề về bảo mật mà chúng ảnh hưởng tới sản phẩm.  Nêu các giấy chứng nhận bảo mật và 
riêng tư cần phải được đáp ứng.> }

\subsection{Các đặc điểm chất lượng phần mềm}
\textit{<Xác định các đặc điểm chất lượng của sản phẩm mà chúng là quan trọng hoặc với nhà phát triển 
hoặc với khách hàng. Một số đặc điểm chất lượng được quan tâm là: tính thích ứng, tính sẵn có, 
tính chính xác, tính linh hoạt,  tính thao tác giữa các phần, tính có thể bảo trì, tính khả chuyển, tính 
tin cậy, tính có thể tái sử dụng, tính có thể kiểm thử, tính dễ sử dụng.  Viết các đặc điểm này cụ 
thể, lượng hóa và có thể kiểm tra khi cần thiết.> }

\subsection{Các quy tắc nghiệp vụ}
\textit{<Liệt kê các nguyên tắc vận hành sản phẩm, chẳng hạn như cá nhân nào hay vai trò nào có thể 
thực hiện các chức năng nào trong các tình huống cụ thể. Bản thân chúng không phải là các yêu 
cầu chức năng nhưng chúng có thể đưa đến các yêu cầu chức năng cụ thể phải tuân theo các 
luật.> }
